\section{Many factors and birational transports}
\label{sec:transports}

We next turn the quadratic cost bound into a birational obstruction.
The argument uses many factors, but their number is fixed before the order
of the covering group tends to infinity.  This order of choices will be
important in the application. Rational evaluation on diagonal products
also appears in algebraic formulations of Lie's superposition theorem;
compare \cite[\S\S3,6--7 and Theorem~7.1]{BlazquezSanzMorales2010}.
We prove the rational evaluation needed here directly from a generically finite
forgetting map of degree one.

\begin{proposition}\label{prop:many-factor-contradiction}
There is no sequence of finite cyclic covers
\[
 \pi_\nu:Y_\nu\longrightarrow X_\nu
\]
of normal projective fourfolds, generically torsors for cyclic groups
\(G_\nu\) of orders \(r_\nu\to\infty\), with the following properties.
The group \(\Bir(X_\nu)\) is countable, and there are ample rational
Cartier divisors \(L_\nu\) on \(X_\nu\) and constants \(c,C>0\),
independent of \(\nu\), such that at very general points
\[
 \varepsilon(L_\nu)\ge c,
 \qquad
 \varepsilon(\pi_\nu^*L_\nu)\ge c r_\nu^{1/4}.
\]
Moreover, for each fixed integer \(t\ge2\), put
\[
 Z_{\nu,t}=Y_\nu^t/G_{\nu,\mathrm{diag}},\qquad
 \theta_{\nu,t}:Z_{\nu,t}\longrightarrow X_\nu^t,
 \qquad
 P_{\nu,t}=\theta_{\nu,t}^*L_\nu^{\boxplus t}.
\]
Here \(L^{\boxplus t}=\sum_{i=1}^t\operatorname{pr}_i^*L\).
The further required bound is
\[
 \varepsilon(P_{\nu,t})\le Ct^2
\]
at a very general point, for all sufficiently large \(\nu\), with this
threshold allowed to depend on \(t\).
\end{proposition}

\begin{proof}
Choose an integer \(T\), to be made sufficiently large in terms of
\(c,C\).  Passing to a tail of the sequence, we may use all the upper
bounds with \(2\le t\le T\) simultaneously.  We suppress \(\nu\)
from the notation until it is necessary to vary the cover.  Write
\(L_Y=\pi^*L\).

Since \(G\) is abelian, the finite map \(\theta_t\) is the quotient
by \(G^t/G_{\mathrm{diag}}\).  Indeed the quotient by the full product
group has function field \(\C(X^t)\), by the generic torsor property,
and the resulting finite birational map to the normal variety \(X^t\)
is an isomorphism.  In particular, \(P_t\) is ample.

\medskip\noindent
\emph{Curve families and compatible leaves.}
Set
\[
 B=CT^2+1.
\]
On each \(Z_t\), choose a marked covering family of integral curves
whose degree divided by multiplicity at the mark is at most \(B\).
Such a family is obtained by choosing a curve at a geometric generic
point witnessing a strict bound below \(Ct^2+1\), and spreading the
curve and its mark.  We use the geometric generic conventions of
Lemma~\ref{lem:generic-chain}: finite data can be defined over a
countable field, and the relevant generic marked families can be
realized over \(\C\).

Enlarge these lists by all deck translates, all permutations of the
slots, and all nonconstant image families under ordered projections
\(Z_t\to Z_s\), for \(2\le s<t\le T\).  Close under these
operations.  The lists remain finite, although their sizes need not be
bounded uniformly in the covering order.

Projection does not increase the marked curve cost.  To see this, let
\(D\to D'\) be the nonconstant map of an integral curve to its reduced
image, let \(e\) be its degree, and let \(z\) and \(z'\) be the marks.
A general local hyperplane through \(z'\), pulled to the normalizations
of the curves, has total order
\(e\mult_{z'}D'\) over \(z'\).  The contribution of the branches
through \(z\) is at least \(\mult_zD\).  Hence
\[
 e\mult_{z'}D'\ge\mult_zD.
\]
Since \(P_t\) minus the pullback of \(P_s\) is nef, we obtain
\[
 \frac{P_s\cdot D'}{\mult_{z'}D'}
 \le
 \frac{P_t\cdot D}{\mult_zD}.
\]
After shrinking parameter spaces we may use reduced integral geometric
generic image curves.  Their projected mark maps remain dominant.

Let \(H_t\) be a geometric generic leaf furnished by
Lemma~\ref{lem:generic-chain}, and put \(h_t=\dim H_t\).
Lemma~\ref{lem:chain-order} gives
\begin{equation}\label{eq:transport-chain-degree}
  h_t>0,\qquad P_t^{h_t}\cdot H_t\le(h_tB)^{h_t}.
\end{equation}
Deck invariance and Corollary~\ref{cor:chain-functoriality} imply that
\(W_t=\theta_t(H_t)\) is the geometric generic fibre closure of a
rational quotient of \(X^t\).  We write
\[
 \xi_t:X^t\dashrightarrow S_t
\]
for such a quotient, with geometrically integral generic fibre.
The leaves and their quotients are equivariant under permutations.

The leaves are also compatible with coordinate projections: an upper
leaf projects into the lower leaf through the projected tuple.  Here is
the genericity justification.  Retain the parameter space of a curve
family when forming its projected image family.  Every parameter and
endpoint occurring in a generic upper chain is generic in its original
parameter space and incidence over the initial field of definition.
Its projected endpoint is generic in the image curve.  For this lower
family, the equality of leaf names at a mark and a generic endpoint is
an identity on the generic incidence.  Consequently it applies to the
projected upper step, regardless of the additional upper history;
a stationary projection gives the same equality immediately.  A path
reaching a generic point of the upper leaf therefore has constant lower
name.  This proves the asserted inclusion.  Whenever different slot
sizes are compared below, we take leaves through a generic tuple and
its projections and extend their name fields to a common algebraically
closed field.

\medskip\noindent
\emph{Finite projections to each slot.}
For \(r\) sufficiently large in terms of \(T,c,C\), the map from
\(H_t\) to each single \(X\)-slot is generically finite onto its
image.  Fix one slot, let its image be \(I\), and write
\[
 \ell=\dim I,\qquad d=h_t-\ell.
\]
The image contains an ambient generic point \(x\) of \(X\), which
is also generic in \(I\) over the leaf name.  The Seshadri lower bound
implies
\begin{equation}\label{eq:transport-image-degree}
 L^\ell\cdot I\ge c^\ell,
\end{equation}
where the assertion is immediate for \(\ell=0\).

Suppose \(d>0\).  Over this generic point \(x\), the geometric fibre
of \(Z_t\to X\) is \(Y^{t-1}\): choose a lift of \(x\) in its
torsor to identify the first coordinate.  The restriction of \(P_t\)
is the external sum of the remaining copies of \(L_Y\).  Every one
of the marked coordinates is generic on \(Y\).  Projecting any curve
through that mark to a nonconstant slot, with the preceding
degree--multiplicity comparison, shows that the Seshadri constant of
this restricted polarization is at least \(cr^{1/4}\).

Let \(J\) be a \(d\)-dimensional component of the geometric generic
fibre of \(H_t\to I\) containing the chosen generic point.  Then
\begin{equation}\label{eq:transport-fibre-degree}
 P_t^d\cdot J\ge (cr^{1/4})^d.
\end{equation}
We recall the elementary degree bound being used twice here.  If an
ample rational divisor has Seshadri constant at least \(a\) at a
smooth point of an \(e\)-dimensional subvariety, blow up that point
and restrict the nef class consisting of the pulled-back divisor minus
\(a\) times the exceptional divisor.  Its top intersection on the
strict transform gives degree at least \(a^e\), since the point has
multiplicity one.  A limiting argument permits a Seshadri bound given
as a supremum.  Our image and fibre points are smooth generic points
of their subvarieties.  The ambient very general bounds continue to
hold on the geometric generic data after the field extensions in use.

The pushforward of \(P_t^d\cdot H_t\) to \(I\) has coefficient
equal to the total degree on the geometric generic fibre, at least
the contribution of \(J\).  Since
\(P_t-\operatorname{pr}^*L\) is nef, equations
\eqref{eq:transport-image-degree} and
\eqref{eq:transport-fibre-degree} give
\[
 P_t^{h_t}\cdot H_t
 \ge (\operatorname{pr}^*L)^\ell P_t^d\cdot H_t
 \ge c^\ell(cr^{1/4})^d.
\]
For fixed \(T\), this contradicts
\eqref{eq:transport-chain-degree} as \(r\to\infty\), because
\(h_t\le4T\).  There are only finitely many slot sizes and slots to
consider.  This proves the assertion simultaneously for all of them
on a tail of the sequence.  In particular,
\begin{equation}\label{eq:transport-ranks}
 1\le h_t\le4.
\end{equation}

\medskip\noindent
\emph{A forgetting map of degree one.}
The same single-slot finiteness holds for \(W_t\).  Let \(D_t\)
denote the degree of its projection to its first-slot image.  The
degree of \(H_t\) over that image contains \(D_t\) as a factor, so
\eqref{eq:transport-chain-degree} and
\eqref{eq:transport-image-degree} imply
\begin{equation}\label{eq:transport-degree-bound}
 D_t\le \frac{P_t^{h_t}\cdot H_t}{c^{h_t}}
       \le K T^8
\end{equation}
where one may take
\[
 K=\max_{1\le h\le4}\left(\frac{h(C+1)}c\right)^h.
\]
Indeed \(h_t\le4\), \(B\le(C+1)T^2\), and \(T^{2h_t}\le T^8\).
Thus \(K\) depends only on \(c,C\), before \(T\) is selected.

Projection compatibility and finite first-slot projection give
\(h_t\le h_{t-1}\).  If these ranks are equal, the image of \(W_t\)
under forgetting its last slot equals the lower leaf: it is contained
in that integral leaf and has the same dimension.  The first-slot
images therefore agree as well.  The degree formula for the two
generically finite maps is
\begin{equation}\label{eq:transport-integer-ratio}
 \frac{D_t}{D_{t-1}}
 =\deg(W_t\longrightarrow W_{t-1})\in\N.
\end{equation}
All degrees are taken after the common geometric field extension
described above, which preserves them.

By \eqref{eq:transport-ranks}, among the \(T-1\) ranks there is a
constant stretch of length at least \((T-1)/4\).  Choose \(T\),
depending only on \(c,C\), so that
\[
 \frac{T-1}{4}-1>\log_2(KT^8).
\]
If every adjacent ratio on that stretch were
larger than one, \eqref{eq:transport-integer-ratio} would force the
degrees to grow by a factor at least two at each step, contradicting
\eqref{eq:transport-degree-bound}.  Hence for some \(k\ge3\),
\begin{equation}\label{eq:transport-birational-omission}
 h_k=h_{k-1}=h>0,
 \qquad W_k\longrightarrow W_{k-1}
 \quad\hbox{is birational}.
\end{equation}
Permutation symmetry gives the same conclusion for every omission of
one slot, with the corresponding lower leaf.  In particular, no
finite correspondence of degree greater than one remains in this
step.

\medskip\noindent
\emph{Rational transport of tangent directions.}
At a generic tuple of \(X^k\), the tangent space along the leaf
\(W_k\) is a rational rank-\(h\) distribution
\(\mathcal D_k=\ker(d\xi_k)\).  It projects injectively to each
single-slot tangent space and isomorphically to the lower distribution
for every omission.  Indeed the corresponding maps of geometric
generic leaves are generically finite, and in characteristic zero
their differentials have rank \(h\) at the generic point.  By the
equal-rank inclusion already proved, their images are exactly the
lower tangent spaces.

Write \(A_i\subset T_{X,x_i}\) for the image in slot \(i\).
This plane depends rationally only on \(x_i\).  To justify the
descent, trivialize \(T_X\) rationally and use a Grassmann chart.
The coordinates of \(A_i\) belong to the function field of each
\((k-1)\)-tuple containing slot \(i\), because they are obtained
from that lower distribution.  The intersection of these coordinate
subfields is the function field of slot \(i\).  Similarly, the graph
of the isomorphism from \(A_i\) to \(A_j\), obtained by lifting to
\(\mathcal D_k\) and projecting to slot \(j\), depends only on
\((x_i,x_j)\).

For completeness, the coordinate-field intersection can be checked
by replacing an omitted coordinate by a fresh independent generic
coordinate.  A function pulled from the other coordinates is unchanged
under that replacement.  Repeating for every omitted coordinate, and
then specializing the fresh coordinates to general constants, shows
that the function belongs to the field of the common coordinates.
This applies to the rational chart coordinates of the planes and
graphs.  It applies also when \(k=3\).

Permutation symmetry thus gives a common rational plane distribution
\(A(x)\) and isomorphisms
\[
 R(x,y):A(x)\xrightarrow{\sim}A(y).
\]
Lifting a vector through the same \(h\)-plane and using three slots
gives the rational cocycle identity
\begin{equation}\label{eq:transport-cocycle}
 R(y,z)R(x,y)=R(x,z).
\end{equation}
Choose a general complex reference point \(a\) so that
\(R(a,-)\) is generically defined and invertible and the specialized
cocycle identity holds.  Transporting a basis of \(A(a)\) gives
rational vector fields \(v_1,\ldots,v_h\) on \(X\), independent at
a generic point.  Equation~\eqref{eq:transport-cocycle} says exactly
that the simultaneous vector fields
\[
 (v_j,\ldots,v_j),\qquad 1\le j\le h,
\]
span \(\mathcal D_k\).  No commutativity assertion about the
\(v_j\) is needed.

\medskip\noindent
\emph{Rational evaluation and local flows.}
Equal ranks were enough to construct the tangent transports. We now
use the stronger degree-one assertion: it turns evaluation of the
missing coordinate into a rational function rather than a finite
correspondence.
Consider the rational map to its image closure
\[
 \Lambda=(\operatorname{pr}_{<k},\xi_k):X^k\dashrightarrow D.
\]
It is birational.  Its restriction to a geometric generic leaf of
\(\xi_k\) is the forgetting map in
\eqref{eq:transport-birational-omission}, and hence has degree one.
More explicitly, let
\[
 F_0\subset E\subset F
\]
be the function fields of \(S_k,D,X^k\), included by pullback.
The field \(E\) is generated over \(F_0\) by the first
\((k-1)\)-slot functions. Since the geometric generic forgetting map
is generically finite, \(F\) and \(E\) have the same transcendence
degree over \(F_0\); hence \(F/E\) is finite. Geometric integrality
of the generic fibre makes \(F/F_0\) regular.  Tensoring with an algebraic closure
\(\overline F_0\) gives fields: each tensor is the union of the
fields obtained by finite algebraic base extension.  Their fraction
fields are the function fields of the geometric generic leaf and its
projected image.  The degree-one geometric projection therefore gives
\([F:E]=1\).  This also rules out an extra generic component or an
unseen algebraic degree in \(\Lambda\).

Let \(\mathrm{ev}\) be the last coordinate of its rational inverse.
Choose a general complex tuple \((u,y_0)\), with
\(u\in X^{k-1}\), in smooth open sets on which the vector fields,
the quotient, and the inverse evaluation are regular.  For small
\(z=(z_1,\ldots,z_h)\in\C^h\), let \(E_z\) be the ordered
composition of the local time-\(z_j\) flows of \(v_j\).  Let
\(E_z^-\) be the reversed composition of the negative-time flows.
These flows exist holomorphically in both the initial point and the
times on sufficiently small neighbourhoods of all chosen coordinates.
One may obtain this directly by Picard iteration in smooth coordinate
polydiscs: bounded coefficients and derivatives make the integral
operator a uniform contraction for small time.  Uniqueness gives the
opposite-time inverse, and reversing the order gives the inverse of
the composition.

The simultaneous flows preserve \(\xi_k\).  Put
\(\lambda(z)=E_z(u)\), acting in all the reference slots.  For
\(y\) near \(y_0\) and sufficiently small \(z\), birational
evaluation therefore gives
\begin{equation}\label{eq:transport-flow-evaluation}
 E_z(y)=\mathrm{ev}\bigl(\lambda(z),\xi_k(u,y)\bigr).
\end{equation}
For each fixed such \(z\), the right side is rational in \(y\).
The equality on an analytic open shows that its input lies in \(D\)
and meets the domain of evaluation; an analytic open is Zariski dense.
Thus it defines a rational self-map of \(X\), regular at \(y_0\)
after shrinking the neighbourhoods and times.

The same argument gives
\[
 E_z^-(y)=\mathrm{ev}\bigl(E_z^-(u),\xi_k(u,y)\bigr)
\]
on a neighbourhood of \(y_0\).  The resulting rational maps are mutually inverse
because they are inverse on a smaller analytic neighbourhood.  They
are therefore birational self-maps of \(X\).  Finally,
\(z\mapsto E_z(y_0)\) has nonzero derivative, since the fields are
independent and \(h>0\).  Its image is uncountable, so the birational
maps, all regular at \(y_0\), are uncountably many distinct maps.
This contradicts the assumed countability of \(\Bir(X)\).
\end{proof}

\begin{lemma}\label{lem:birational-countability}
Let \(X\) be a normal projective terminal variety with
\(K_X\sim_\Q0\).  If a smooth projective resolution of \(X\) has
irregularity zero, then \(\Bir(X)\) is countable.
\end{lemma}

\begin{proof}
Let \(p:W\to X\) be a smooth projective resolution. Since
\(H^1(W,\OO_W)=0\), the exponential sequence embeds the analytic
Picard group in \(H^2(W,\Z)\). By GAGA the analytic and algebraic
Picard groups agree, so \(\Pic(W)\) is countable. Divisor pushforward
induces a surjection \(\Pic(W)\to\Cl(X)\): the strict transform of
every Weil divisor on \(X\) is Cartier on \(W\). Thus \(\Cl(X)\)
is countable.

Every birational self-map of \(X\) is an isomorphism in codimension
one.  To see this, take a smooth common resolution of its graph, with
projections \(p_1,p_2\) to \(X\).  Choose a nonzero trivialization
\(\eta\) of a sufficiently divisible pluricanonical bundle of
\(X\).  Its two pullbacks are nonzero regular pluricanonical forms
on the common resolution.  That section space has dimension one:
canonicality gives pullback regularity, and restriction to a big open
of \(X\), followed by divisorial extension, identifies the sections
with those of the trivial pluricanonical bundle on \(X\).
Consequently the two pullbacks are proportional.  Terminality says
that their zero divisors have support exactly the exceptional prime
divisors of \(p_1\) and \(p_2\), respectively.  These supports must
coincide.  A prime divisor contracted by the self-map or its inverse
would belong to only one of the two supports, which is impossible.

Fix a very ample divisor \(A\) on \(X\).  Group the birational
self-maps according to the Weil linear equivalence class of the strict
transform of \(A\) back to the source.  Within one class, the
complete spaces of sections are identified by the codimension-one
isomorphisms.  The two maps to the embedding of \(X\) by \(|A|\)
therefore differ by a projective linear automorphism preserving that
embedded variety.  It is enough to prove that this stabilizer has
finite effective image on \(X\).

Let $S$ be the identity component of the stabilizer in the projective
linear group of the embedding. Its action on $X$ induces a rational
action on the smooth projective resolution $W$. The effective image is
a connected linear algebraic group: the kernel of the action is a closed
normal subgroup, and a quotient of a linear algebraic group is linear.
As proved above, $W$ has a nonzero positive pluricanonical section.
Matsumura's obstruction therefore says that this effective image is
trivial \cite[Corollary~1]{Matsumura1963}. More precisely, that corollary
excludes positive-dimensional linear algebraic subgroups of $\Bir(W)$
for a smooth complete variety with a nonzero pluricanonical system.
It applies here to the rational action, without requiring that every
birational self-map of $W$ be regular.

The identity component of the stabilizer consequently acts trivially.
An algebraic group has only finitely many components, so the effective
stabilizer is finite.  There are countably many possible divisor
classes and finitely many birational maps in each such class.  Hence
\(\Bir(X)\) is countable.
\end{proof}
